\subsection{Polynomial iid learning for the same nonlinear prior}
\label{sec:near-iid}

For the near-identity family, the exponential online lower bound is
accompanied by a polynomial iid upper bound for exactly the same prior,
activation, target, and width. This does not require a new restriction
on the learner's parameter support.

\begin{corollary}[An iid/online separation at every width]
\label{cor:iid}
Fix the hypotheses of Theorem~\ref{thm:near} or the larger perturbation
allowed by Corollary~\ref{cor:near-half}, including $k,L,\phi$, and
any finite width $N\ge1$. Let $\mu$ be any distribution on the input
cube. For $\varepsilon,\eta\in(0,1)$, suppose
\begin{equation}\label{eq:iid-sample}
 n\ge\frac6\varepsilon
       \left([1+d(d+1)]\log2+\log(1/\eta)\right).
\end{equation}
For every sufficiently small positive $\sigma$, with probability at least
$1-\eta$ over $n$ iid inputs labeled by $f$, the posterior gives the
correct label probability strictly greater than $2/3$ on a set of
$\mu$-measure at least $1-\varepsilon$. Its posterior-majority test
error and its mean test error under fresh posterior sampling are also
at most $\varepsilon$.

The bound on $n$ does not depend on learner width or depth. The
sufficiently-small-$\sigma$ threshold may depend on $k,L,N,\phi$ and
$\varepsilon$, but can be chosen independently of $\mu$, $n$, and the
particular target-consistent history. The activation remains a
size- and depth-dependent near-identity family.
\end{corollary}
\begin{proof}
Let $\mathcal H_d$ be the class of affine-threshold sign functions on
the cube, and set
\[
 U_*=L\left(\frac{dQ_*}{4\alpha^{2(L-1)}}\right)^{1/L}.
\]
Every weak-global-margin parameter with energy at most $U_*$ has an
affine sign function. Indeed, the linearization bound
\eqref{eq:linear-error} and the calculation \eqref{eq:delta}, now with
$t=M$, bound the error by
$1/[128(d+1)Q_*]<1/2$. Under the larger perturbation in
Corollary~\ref{cor:near-half}, the same calculation as
\eqref{eq:delta-half}, with $t=M$, instead bounds it by
$2^{k+1}/(128Q_*)\le1/128<1/2$. In either case its affine linearization
has the same sign as its nonlinear output everywhere.

The union $B$ of weak-margin cells of nonaffine Boolean functions is
closed, since the input domain is finite. It contains no point of
energy at most $U_*$. Coercivity implies $\inf_B\kappa_{d,N}>U_*$
unless $B$ is empty. On the other hand, the scalar-path construction
with $A_M$ gives the strict full-target cell
\[
 U_f=\{\theta:f(\xi)Z_\theta(\xi)>1/2\text{ for every }\xi\}
\]
energy infimum at most $U_*$. The Gaussian bounds used in
Lemma~\ref{lem:gaussian} show that
\[
 R_\sigma:=\frac{\Pp_\sigma(B)}{\Pp_\sigma(U_f)}\longrightarrow0.
\]
For every history $H$ labeled by $f$, its conditioning event $E_H$
contains $U_f$, while its nonaffine part is contained in $B$. Thus
\begin{equation}\label{eq:uniform-affine}
 \Pi_\sigma(\mathcal H_d^c\mid H)
 \le\frac{\Pp_\sigma(B)}{\Pp_\sigma(U_f)}=R_\sigma.
\end{equation}
This is one bound for all such histories, without a union bound over
histories or any dependence on the sampling distribution.

A hyperplane-arrangement count gives
\[
 |\mathcal H_d|\le2\sum_{j=0}^{d}\binom{2^d-1}{j}
                 \le2^{1+d(d+1)}.
\]
For any $\rho>0$, a fixed classifier with $\mu$-error greater than
$\rho$ remains consistent with $n$ iid labels with probability at most
$e^{-n\rho}$. A union bound implies that, with probability at least
$1-\eta$, every consistent affine classifier has test error at most
$\rho$, provided
\[
 n\ge\frac{[1+d(d+1)]\log2+\log(1/\eta)}\rho.
\]
Choose $\rho=\varepsilon/6$ and choose $\sigma$ small enough that
$R_\sigma\le\varepsilon/6$. On this sample event, the mean posterior
error at a fresh input is at most
$\rho+R_\sigma\le\varepsilon/3$.
Write $q(\xi)$ for the posterior probability of the wrong label.
Posterior-majority errors are contained in $\{q\ge1/2\}$, and failures
of strict $2/3$ confidence are exactly $\{q\ge1/3\}$.
Markov's inequality bounds their $\mu$-measures by
$2\varepsilon/3$ and $\varepsilon$, respectively, proving the result.
\end{proof}

Thus the same nonlinear learner has an
$O((d^2+\log(1/\eta))/\varepsilon)$ iid sample guarantee and at least
$2^k-k-1$ designated online errors, or $2^{k-1}$ under the larger
perturbation allowance, with $d=2k+1$, at every fixed width.
For fixed $k,L,N$, one sufficiently small positive scale satisfies both
claims. The contrast is between different ways of choosing examples;
it does not require changing the prior during learning.
